\section{Flowering: the per-flower Markov chain}\label{sec:flowering}
\stagebox{\statuspaused. Only the 2D Python prototype exists.}{\src{prototype-python/automark_bloom/petal_markov.py} (prototype); 3D core not started}

Flowering takes over once the bud structure is fixed. Each growth step depends
probabilistically on the previous state, a semi-random walk, and each flower
draws from its own stream (key $1000+i$ for flower $i$,
\cref{sec:seeds}). Flowers on one plant therefore share their bud and their
targets but take different paths: siblings, not clones.

Andres has put the 3D flowering stage on hold. This section records the
prototype that exists and the design agreed so far, and will be rewritten when
work resumes.

\subsection{The prototype (2D, discrete)}
\paragraph{State.} Each petal carries three bin indices: length $\lambda$ (14 bins,
0.08 to 0.72 flower radii), width $\omega$ (12 bins, 0.06 to 0.60) and curl
$\gamma$ (13 bins, $-20^\circ$ to $100^\circ$; positive cups inward, negative
bends outward). Every petal starts as a closed bud: shortest, narrowest and
fully curled. The walk runs 45 steps.

\paragraph{Targets.} Each petal has a target bin for each variable, set by its
place in the bud: outer petals end long, wide and flat, inner petals short and
cupped. For a petal whose outer edge sits at relative radius $\pi_1\in[0,1]$,
the target length is $0.22 + 0.48\,\pi_1$ and the target curl is
$80^\circ - 95^\circ\pi_1$; the target width comes from the petal's angular span
in the bud grid.

\paragraph{Moves.} At every step each index $i$ moves by $+1$, $0$ or $-1$.
With $d = i^\star - i$ the (fractional) distance to the target bin and
$\sigma(z) = 1/(1+e^{-z})$,
\[
  \Pr(+1) = g_\uparrow\,\sigma(1.5\,d), \qquad
  \Pr(-1) = g_\downarrow\,\sigma(-1.5\,d), \qquad
  \Pr(0) = 1 - \Pr(+1) - \Pr(-1),
\]
so the walk drifts toward the target but keeps jittering around it. Indices are
clipped to their bins.
\begin{center}\small
\begin{tabular}{@{}lcc@{}}
\toprule
Variable & $g_\uparrow$ & $g_\downarrow$ \\
\midrule
length & 0.55 & 0.20 \\
width & $0.50\,e^{-1.5P}$ & 0.20 \\
curl & 0.35 & 0.45 \\
\bottomrule
\end{tabular}
\end{center}

\paragraph{Packing pressure.} $P$ is a petal's fractional angular overlap with
neighbours within two rings, computed from the current state. Crowding
suppresses widening (the $e^{-1.5P}$ factor) and raises the curl target by
$30^\circ\min(P,1)$, so crowded petals bend instead of spreading. Because $P$
couples petals, the Markov property holds for the joint state of all petals,
not for each petal on its own.

\subsection{The 3D design agreed so far}
\begin{itemize}[leftmargin=*]
\item \textbf{Same object throughout.} Flowering starts from each petal's
  packed sheet in the bud (\cref{sec:sheets,sec:packing}) and reshapes those
  same control points, so the opening animation starts from the packed bud.
\item \textbf{Continuous state.} Length, width and curl change continuously
  rather than in bins, while the steps stay random.
\item \textbf{The chain moves targets; the energy moves petals.} The Markov
  chain changes each petal's targets (for example rest lengths and preferred
  curvature), and the sheets follow by relaxing the same differentiable energy,
  so petals keep their thickness and never pass through each other while
  opening.
\item \textbf{Frames.} Each step becomes a frame in the flower JSON with the
  same topology (\cref{sec:output}).
\item \textbf{Ideas to fold in} from the design notes: phyllotactic spiral
  angle, petals growing larger outward, and bending to pack against neighbours.
\end{itemize}
